\section{Introduction}\label{sec:intro}

Many organisations still run business-critical applications on SAP ASE, formerly Sybase, and a common
modernisation path is migration to Microsoft SQL Server. Practitioner guidance describes this path as a
sequence of assessment, planning, schema and code conversion, data movement, testing, cut-over, and
tuning~\cite{manda2026}. Vendor tools such as SQL Server Migration Assistant (SSMA)~\cite{ssma} convert
much of the schema and code, but stored procedures, triggers, and embedded SQL still need manual work, and
most of the effort and risk of a migration sits there.

LLMs can help with that work, but their output is hard to trust. ASE and SQL Server both descend from the
same Transact-SQL code base, so most ASE code is close to valid SQL Server code. Some of the remaining
differences cause compile errors: the ASE \texttt{*=} outer join, for example, is rejected by SQL Server.
Others do not. Under ASE's default settings, \texttt{WHERE col = NULL} matches NULL values, while a SQL
Server procedure created with \texttt{ANSI\_NULLS ON} returns no rows for the same predicate. Concatenating
a NULL with a string gives the string in ASE and NULL in SQL Server. Moving a condition on the inner table
of an ASE outer join into the \texttt{WHERE} clause turns it into an inner join. A migration process that
only checks whether translated code compiles will accept all of these.

This paper studies how well locally hosted LLMs translate ASE code to SQL Server and which checks are
needed before their output can be used. We run the models locally because migration source code often
contains business rules that organisations do not send to external services. Since no ASE server was
available, all results measure agreement with reviewed reference translations; Section~\ref{sec:taxonomy}
explains how the references were built and where they depend on documented ASE behaviour. The paper
contributes:
\begin{enumerate}
\item A taxonomy of \NumCategories{} categories of ASE-to-SQL Server differences, each labelled by how an
unconverted construct shows up, and an open corpus of \NumDevCases{} development cases and \NumHeldout{}
held-out procedures with reference translations (Section~\ref{sec:taxonomy}).
\item VeriMig, a translation framework with four verification gates and a repair loop that never sees the
reference outcomes (Section~\ref{sec:framework}).
\item An evaluation of a rule baseline, plain and checklist prompting, and the repair loop across five
model configurations, with an analysis of the failures that pass compilation
(Sections~\ref{sec:setup}--\ref{sec:results}).
\item A fault-injection study of the data-reconciliation checks used at cut-over, including the encoding
needed for row hashes to be reliable.
\end{enumerate}
We see the benchmark and the failure analysis as the main contributions. The framework combines known
techniques; its value is in how they are arranged and measured for this migration.
